\chapter{Development Microeconomics}\label{ch:b3:development-microeconomics}

The chapters before this one asked why whole countries are poor. This one asks what
changes the lives of poor households: credit, cash, schooling, health, secure title
to a home. For each question there are now randomised experiments, often several,
and the interest has moved from whether a programme works somewhere to how much it
does, for whom, and whether the answer travels. The chapter starts with how poverty
is measured, because what counts as success depends on it.

\begin{openingchart}{The general effect of expanding access to microcredit on six
household outcomes, pooled over seven randomised experiments (Bosnia, Ethiopia, India,
Mexico, Mongolia, Morocco, the Philippines) by a Bayesian hierarchical model: posterior mean and central 95~\%
interval, dollars at purchasing power parity per two weeks (2010).}{Meager (2016),
arXiv:1506.06669v3, Appendix A (joint model)}
\begin{tikzpicture}
\begin{axis}[omscatter, width=0.66\linewidth, height=0.42\linewidth, xmin=-10,
  xmax=46, ymin=0.4, ymax=6.6, ytick={1,...,6},
  yticklabels={Temptation goods, Consumer durables, Consumption, Expenditures, Revenues, Profit},
  xlabel={Dollars PPP per two weeks},
  extra x ticks={0}, extra x tick labels={}, extra x tick style={grid=major, grid style={black!60}}]
  \addplot+[mark=*, mark size=2pt, error bars/.cd, x dir=both, x explicit]
    table[x=tau, y=pos, x error minus=el, x error plus=eh] {figdata/bachelor-3/34-development-microeconomics/out/microcredit.csv};
\end{axis}
\end{tikzpicture}
\end{openingchart}

% figdata: microcredit
% ledger: b3dm:meager.tab, b3dm:meager
Every interval includes zero. The largest effect is on business revenues, 14.45 dollars per two weeks, and on profit it is
6.81 --- 7~\% of the control group's mean of 94.81. On consumption it is 3.44, about 1~\% of the mean. The effect of
\omterm{def:b3:development-microeconomics:microcredit}{microcredit} is likely positive and small, and a negative effect cannot be ruled out.
Pooling all the data as if the seven sites were one gives ``significant'' results for 2 of the six
outcomes. The six studies published together describe the same pattern: modestly positive, not
transformative. % ledger: b3dm:bkz

\begin{recall}
The poverty line and the poverty rate are the grades 6--9 volume's; randomised trials
and whether a result travels are the Year 2 volume's; \omterm{def:b3:instrumental-variables:instrument}{instrumental variables} and
regression discontinuity are \cref{ch:b3:instrumental-variables,ch:b3:regression-discontinuity}'s;
\omterm{def:b3:credible-empirical-work:meta-analysis}{meta-analysis} is \cref{ch:b3:credible-empirical-work}'s; \omterm{def:b3:human-capital:human-capital}{human capital} is
\cref{ch:b3:human-capital}'s.
\end{recall}

\begin{notation}[Symbols of this chapter]\label{not:b3:development-microeconomics:symbols}
$z$ is the poverty line; $y_i$ income or consumption of person $i$; $N$ the population; $P_\alpha$ the
Foster--Greer--Thorbecke measure of order $\alpha$; $\kappa(\cdot)$ work capacity; $T$ a transfer.
\end{notation}

\section{Measuring poverty in depth}\label{sec:b3:development-microeconomics:measuring}

The headcount --- the share of people below the line --- does not change when a
poor person becomes poorer, or when one just below the line is helped while one far
below it is not. Two refinements fix this.

\begin{definition}[Poverty gap index; squared poverty gap]\label{def:b3:development-microeconomics:fgt}
\index{poverty gap index}\index{squared poverty gap}
For a line $z$, the Foster--Greer--Thorbecke measures are $P_\alpha = \frac1N\sum_{y_i < z}\left(\frac{z - y_i}{z}\right)^{\alpha}$. $P_0$ is the headcount; $P_1$, the
\emph{poverty gap index}, is the average shortfall from the line as a share of it,
counting the non-poor as zero; $P_2$, the \emph{squared poverty gap}, weights each
shortfall by itself. % ledger: b3dm:fgt
\end{definition}

\begin{proposition}[What the measures respond to]\label{prop:b3:development-microeconomics:fgt}
$P_1$ and $P_2$ rise when a poor person's income falls (monotonicity). $P_2$ also rises
when a sum is transferred from a poor person to a less poor person who stays poor
(transfer sensitivity); $P_1$ does not change.
\end{proposition}

\begin{proof}
For $y_i < z$, $\partial P_\alpha/\partial y_i = -\frac{\alpha}{Nz}((z - y_i)/z)^{\alpha - 1} < 0$ for $\alpha \ge 1$: monotonicity. A transfer $t$ from $i$ to $j$, with $y_i < y_j$ and
both poor before and after, changes $P_1$ by $\frac{1}{Nz}(t - t) = 0$. For $P_2$, the change is $\frac1N[g(y_i - t) - g(y_i) + g(y_j + t) - g(y_j)]$ with $g(y) = ((z - y)/z)^2$, convex and
decreasing on $[0, z]$; since $y_i - t < y_i \le y_j < y_j + t$, the rise at $i$ exceeds the fall at $j$ by convexity.
\end{proof}

$N z P_1$ is the sum of all shortfalls: the cost of ending poverty with transfers aimed
perfectly, which no real programme achieves.

\begin{omchart}{Headcount ($P_0$), poverty gap ($P_1$) and squared poverty gap ($P_2$) at 3.00 dollars a day
(2021 PPP), eight countries, 2022, in percent. Computed from the percentile groups of
the World Bank's Poverty and Inequality Platform (each group at its mean, so the
headcount moves in steps of one point; consumption, except income for Brazil).}{World
Bank, Poverty and Inequality Platform, world\_100bin}
\begin{tikzpicture}
\begin{axis}[omrecord, xbar, bar width=3pt, width=0.62\linewidth, height=0.5\linewidth,
  xmin=0, xmax=85, ytick={1,...,8},
  yticklabels={Brazil, India, Bangladesh, Indonesia, South Africa, Nigeria, Kenya, Mozambique},
  xlabel={Percent}, legend pos=south east, enlarge y limits=0.08]
  \addplot[fill=omDef!70, draw=omDef] table[x=p0, y=pos] {figdata/bachelor-3/34-development-microeconomics/out/fgt.csv};
  \addplot[fill=omThm!60, draw=omThm] table[x=p1, y=pos] {figdata/bachelor-3/34-development-microeconomics/out/fgt.csv};
  \addplot[fill=black!30, draw=black!60] table[x=p2, y=pos] {figdata/bachelor-3/34-development-microeconomics/out/fgt.csv};
  \legend{$P_0$, $P_1$, $P_2$}
\end{axis}
\end{tikzpicture}
\end{omchart}
% figdata: fgt

Kenya and Nigeria have almost the same headcount, 45 and 42~\%, and almost the same gap, 13.6 and
13.0~\%. Mozambique's headcount is 81~\%, but its gap, 44.7~\%, is more than three times Kenya's: its poor are
much further below the line. India, Bangladesh and Brazil have headcounts near 5--6~\%; Brazil's gap is
twice India's, but Brazil's survey measures income and India's consumption, and
income varies more, so part of the factor is the welfare measure.

\begin{measure}[The FGT poverty measures]
The measures need the whole distribution below the line, which comes from household
surveys of consumption (most poor countries) or income (most of Latin America); the
two are not the same, since income varies more than consumption. They also inherit
the price index that converts the line into local currency.
\end{measure}

\section{Health, nutrition and the poverty trap}\label{sec:b3:development-microeconomics:health}

\begin{definition}[Nutrition-based poverty trap]\label{def:b3:development-microeconomics:poverty-trap}
\index{nutrition-based poverty trap}
In a \emph{nutrition-based poverty trap}, a worker's capacity to work $\kappa$ depends on
what she eats, and so on her income: $y_{t+1} = w\kappa(y_t)$. If capacity rises slowly at very low intake, then
steeply, then slowly again, the poor stay poor because they are too weak to earn more.
\end{definition}

\begin{proposition}[A trap needs an S-shape]\label{prop:b3:development-microeconomics:trap}
If $f(y) = w\kappa(y)$ is increasing, concave and bounded with $f(0) > 0$, there is exactly one fixed point, and it is
stable: no trap. If $f$ is convex and then concave and crosses the 45-degree line three times, at $y_L < y_M < y_H$,
then $y_L$ and $y_H$ are stable and $y_M$ unstable: households starting below $y_M$ converge to $y_L$, those
above to $y_H$.
\end{proposition}

\begin{proof}
With $f$ concave, $f(y) - y$ is concave, positive at 0 and, since $f$ is bounded, negative for large $y$, so it
has one zero $y^*$, where $f$ crosses from above: $f'(y^*) < 1$, stable. With three crossings, $f(y) - y$ changes sign
from positive to negative at $y_L$ and $y_H$ (slopes of $f$ below 1, stable) and from negative to positive at $y_M$
(slope above 1, unstable); iterating $f$ from below $y_M$ moves $y$ towards $y_L$, from above towards $y_H$.
\end{proof}

\begin{model}{The poverty trap}
\assumes Returns to income or assets that are increasing at low levels: a minimum
scale of business, nutrition below working capacity, or lumpy assets.
\predicts A one-time transfer that lifts a household above the unstable point puts
it on a path to the high state; a smaller one is eaten up and the household returns
to the low state.
\tested In Bangladesh, over 21,000 households in 1,309 villages were randomised into a programme that gives
the poorest women livestock and skills; they moved from casual wage labour into
livestock rearing, accumulated assets and left poverty, with the gains sustained
after four and seven years. % ledger: b3dm:bandiera
\breaks A review of the mechanisms --- savings, coordination, nutrition, business
size --- finds such traps rare and largely confined to remote or disadvantaged
areas, with geography the most likely source. % ledger: b3dm:km
\end{model}

Household calorie needs are met at low incomes in most of the world, which is one
reason the nutrition trap is rare: the poor spend extra income partly on better
food, not on more calories.

\begin{record}[Deworming]
School-based deworming, randomly phased into Kenyan schools, cut absenteeism by a quarter
and benefited untreated children in treated and neighbouring schools; it did not
improve test scores. % ledger: b3dm:mk
A pure \omterm{def:b3:credible-empirical-work:replication}{replication} with the original data reproduced most results, but after
coding errors were corrected the effect on nearby schools largely disappeared, and
the combined gain in attendance among treated children fell from 7.5~\% to 3.9~\%, no longer
statistically significant. % ledger: b3dm:aiken
The exchange is the lesson of \cref{ch:b3:credible-empirical-work}: replicate, and
separate what the data say from the specification that produced it.
\end{record}

\section{Schooling and cash transfers}\label{sec:b3:development-microeconomics:cash}

\begin{definition}[Conditional cash transfer; unconditional cash transfer]\label{def:b3:development-microeconomics:cash-transfer}
\index{conditional cash transfer}\index{unconditional cash transfer}
A \emph{conditional cash transfer} pays poor households on condition that they do
something, typically keep their children in school or attend health checks. An
\emph{unconditional cash transfer} pays them without conditions.
\end{definition}

\begin{record}[Progresa, 1997]
Mexico's Progresa paid education grants to poor mothers in rural communities.
Communities were randomly chosen for the first phase, and the evaluation compared
children in them with children in control communities, checking the comparison
before the programme and computing double differences by grade and sex; the
short-run effects on enrolment were extrapolated to lifetime schooling and earnings
to give a rate of return on the grants. % ledger: b3dm:progresa
The design --- phasing a national programme in at random --- became the model for
the evaluations that followed.
\end{record}

\begin{record}[GiveDirectly in Kenya]
Unconditional transfers of 404 or 1{,}525 dollars (PPP) to poor rural households, randomised at the level of
villages and households, raised monthly household consumption from 158 to 193 dollars nine months after the
transfers began, and raised psychological well-being; monthly payments went more to
food, lump sums more to durables. % ledger: b3dm:hs
A later trial gave about 1{,}000 dollars each to over 10,500 households in 653 randomised villages, a shock of over 15~\%
of local GDP; non-recipient households and firms gained too, prices barely rose, and the
local multiplier was estimated at 2.4. % ledger: b3dm:egger
\end{record}

The second study answers an objection to the first: if transfers raised prices,
part of the measured gain would be nominal, and non-recipients would lose. They did
not, at least where local supply could expand.

\section{Credit and microfinance}\label{sec:b3:development-microeconomics:credit}

\begin{definition}[Microcredit; group liability]\label{def:b3:development-microeconomics:microcredit}
\index{microcredit}\index{group liability}
\emph{Microcredit} is lending of small sums, without collateral, to poor households,
usually for a business. Under \emph{group liability} the members of a group of
borrowers are each responsible for the others' loans, which gives them a reason to
choose partners well and to watch one another.
\end{definition}

\omterm{def:b3:development-microeconomics:microcredit}{Group liability} raises the probability of repayment when partners can cover each
other (exercise 7), so lenders can charge less. The opening chart is the result of
seven experiments that expanded access to such loans: the general effect is small
and uncertain. The disagreement across sites comes mostly from the 27~\% of households that already ran a
business, which gain much more; about 60~\% of the variation across sites is pooled, so the sites are
informative about one another. % ledger: b3dm:meager
For a new site, the effect on profit has a 95~\% chance of lying between $-16$ and 40 dollars per two weeks, against an interval of
$-2$ to 17 if the sites are pooled as one: ignoring the heterogeneity overstates the precision. % ledger: b3dm:meager

\section{Property titles}\label{sec:b3:development-microeconomics:titles}

\begin{definition}[Property titling]\label{def:b3:development-microeconomics:titling}
\index{property titling}
\emph{Property titling} gives the occupants of land or a dwelling a registered legal
title to it.
\end{definition}

\begin{record}[Land titles in Buenos Aires]
In 1981 squatters occupied land in a poor suburb of Buenos Aires; a law of 1984 expropriated the owners to
give the occupants title. Some owners accepted compensation and others disputed it
in slow courts, so some squatters received titles and their neighbours did not, for
reasons unrelated to them. Surveys in 2003 and 2007 found that titled families invested more in their
houses, had smaller households and educated their children more --- but not through
more credit. % ledger: b3dm:gs
\end{record}

The result fits \cref{prop:b3:institutions:holdup} of the previous chapter better
than the hope that titles unlock credit: security raised investment directly.

\section{Aid and growth}\label{sec:b3:development-microeconomics:aid}

\begin{omchart}{Net official development assistance received (\% of GNI) and growth of GDP
per person, country averages 1990--2019, 134 countries with at least twenty years of both, with the least-squares line.}{World Bank, World Development Indicators,
DT.ODA.ODAT.GN.ZS and NY.GDP.PCAP.KD.ZG}
\begin{tikzpicture}
\begin{axis}[omscatter, width=0.6\linewidth, height=0.4\linewidth,
  xlabel={Aid, \% of GNI}, ylabel={Growth, \% a year}, xmin=0]
  \addplot table[x=aid, y=growth] {figdata/bachelor-3/34-development-microeconomics/out/aid.csv};
  \addplot[black!60, thick, no marks, sharp plot] table[x=aid, y=growth] {figdata/bachelor-3/34-development-microeconomics/out/aid_fit.csv};
\end{axis}
\end{tikzpicture}
\end{omchart}
% figdata: aid_stats

Across 134 countries the correlation between aid and growth is $-0.16$: each point of GNI received in aid goes
with 0.04 points less growth. The correlation says little: aid goes to countries that are poor,
in crisis or at war, which grow slowly for other reasons.

\begin{debate}[Aid and growth]
\textbf{Aid works where policy is good.} Aid raises growth in developing countries
with good fiscal, monetary and trade policies and has little effect where policies
are poor (\emph{Grade 3}). % ledger: b3dm:bd
\emph{Weakness:} an interaction found in one dataset.

\textbf{The interaction does not survive new data.} Extending the data from 1970--93 to 1970--97 and filling
the gaps, the finding is no longer robust (\emph{Grade 3}). % ledger: b3dm:elr
\emph{Weakness:} a failure to replicate is not a proof of no effect.

\textbf{No robust effect.} After correcting for aid being drawn by poor or strong
growth, there is little robust evidence of a positive or a negative effect, and none that
aid works better with better policies or that some forms work better (\emph{Grade 3}). % ledger: b3dm:rs
\emph{Weakness:} the instruments for aid in cross-country work are themselves
debated.

\textbf{A quasi-experiment.} Since 1987, eligibility for the World Bank's concessional
lending has depended partly on an income threshold; aid falls by about 59~\% of its share of GNI after a
country crosses it. For the 35 countries that crossed from below in 1987--2010, one more point of aid in GNI raises
growth of income per person by about 0.35 points a year (\emph{Grade 2}). % ledger: b3dm:gkxz
\emph{Weakness:} it estimates the effect for countries near the threshold, and
crossing it may change more than aid.
\end{debate}

\section{Whether a result travels}\label{sec:b3:development-microeconomics:travels}

A programme that works in one place may not work in another: the people, the
markets, the prices and the implementer differ. Two tools from this volume help. A
hierarchical model, as in the opening chart, estimates how much the effects differ
across sites and gives an interval for a new site; and the general-equilibrium
effects of a programme --- on prices, wages and non-participants --- must be
measured at the scale at which it would run, as the second Kenyan trial did by
randomising whole villages. A small trial that changes no price can hide a large
programme that does.

\section{Exercises}\label{sec:b3:development-microeconomics:exercises}

\begin{exercise}[$\star$]\label{exo:b3:development-microeconomics:1}
From the opening chart, give the effect on profit as a share of the control mean, and
say which outcome has the only negative point estimate.
\end{exercise}

\begin{exercise}[$\star$]\label{exo:b3:development-microeconomics:2}
Five people have incomes 2, 4, 6, 10 and 12, and the line is 8. Find $P_0$, $P_1$ and $P_2$.
\end{exercise}

\begin{exercise}[$\star$]\label{exo:b3:development-microeconomics:3}
How much would perfectly targeted transfers to the people of exercise 2 cost to end
poverty? Express it with $P_1$.
\end{exercise}

\begin{exercise}[$\star$]\label{exo:b3:development-microeconomics:4}
In an asset-transfer experiment with randomised transfer sizes, what result would
contradict the poverty-trap model?
\end{exercise}

\begin{exercise}[$\star\star$]\label{exo:b3:development-microeconomics:5}
Compute the three FGT measures for Kenya in 2022 at 3.00, 4.20 and 8.30 dollars a day.
\end{exercise}

\begin{exercise}[$\star\star$]\label{exo:b3:development-microeconomics:6}
Three invented studies estimate 10 (standard error 4), 2 (2) and 6 (4). Find the inverse-variance pooled
estimate and its standard error.
\end{exercise}

\begin{exercise}[$\star\star$]\label{exo:b3:development-microeconomics:7}
Two borrowers each borrow 1 and must repay $r$; each project succeeds with probability 0.8, independently,
and returns 3. Under \omterm{def:b3:development-microeconomics:microcredit}{group liability}, a successful borrower repays for a failed
partner. Find the probability that the lender is repaid in full, with and without
\omterm{def:b3:development-microeconomics:microcredit}{group liability}, and the $r$ at which a lender whose funds cost 1 breaks even.
\end{exercise}

\begin{exercise}[$\star\star$]\label{exo:b3:development-microeconomics:8}
A conditional transfer pays 400 dollars a year for each enrolled child; enrolment rises from
70 to 80~\%. Find the cost per extra child-year of schooling.
\end{exercise}

\begin{exercise}[$\star\star$]\label{exo:b3:development-microeconomics:9}
What evidence would show that the measured rise in recipients' consumption after a
transfer is a rise in prices, not in what they consume?
\end{exercise}

\begin{exercise}[$\star\star\star$]\label{exo:b3:development-microeconomics:10}
Prove that $P_2$ is transfer-sensitive and that $P_0$ can fall when a transfer from a very poor to a
nearly poor person makes the second non-poor.
\end{exercise}

\begin{exercise}[$\star\star\star$]\label{exo:b3:development-microeconomics:11}
Seven invented studies of a training programme give effects 12, 3, 25, $-4$, 8, 15, 1 with standard errors 5, 4,
9, 6, 3, 7, 4. Find the fixed-effect estimate, Cochran's $Q$, $\tau^2$ and the random-effects estimate.
\end{exercise}

\begin{exercise}[$\star\star\star$]\label{exo:b3:development-microeconomics:12}
Local supply is $Q = P^{\varepsilon}$, and a transfer raises local spending $PQ$ by 15~\%. Find the rise in prices for $\varepsilon = 0$
and $\varepsilon = 9$. Which fits the Kenyan evidence?
\end{exercise}

\section{Problem: A Cash Transfer for Brenland's Poorest}\label{sec:b3:development-microeconomics:problem}

\begin{problem}[{Weekend problem --- A cash transfer for Brenland's poorest}]\label{pb:b3:development-microeconomics:cash}
The 2025 household survey (\code{load("households")}) gives each household's size, disposable
income and consumption. The poverty line is 60~\% of the median of disposable income per
person, counted over people. An invented programme covers every household below the
line: half, drawn at random, receive one transfer of 3{,}000, 6{,}000 or 9{,}000 dalens (also at random), the
others are controls. Two years later the programme measures productive assets and
consumption per person. % figdata: cash_stats

\medskip\noindent\textbf{Part I --- poverty in Brenland.}
\begin{enumerate}
\item Find the median of disposable income per person and the line.
\item Find $P_0$, $P_1$ and $P_2$ over people.
\item Find the total yearly sum that perfectly targeted transfers would need.
\item How many households fall below the line?
\item Why is the survey's consumption a better welfare measure here than income?
\item Why might the headcount be the wrong target for a programme?
\end{enumerate}

\medskip\noindent\textbf{Part II --- the randomised transfer.}
\begin{enumerate}[resume]
\item How many households are treated and how many are controls?
\item Compare baseline consumption per person in the two groups. Why should it be close?
\item Find mean consumption per person after two years in each group.
\item Estimate the effect of the transfer and its standard error.
\item Give the 95~\% interval.
\item Find the share of each group above the line after two years.
\item Why does the control group's share above the line rise at all?
\end{enumerate}

\medskip\noindent\textbf{Part III --- a threshold?}
\begin{enumerate}[resume]
\item Find mean assets after two years for controls and for each transfer size.
\item For each transfer size, find the share of households whose assets plus transfer
reach 7{,}000 dalens.
\item What happens to the assets of the households with 3{,}000? With 9{,}000?
\item Does the pattern support a poverty trap? What would a proportional effect look
like?
\item Why does randomising the size of the transfer matter for this test?
\item Find the share above the line after two years for each transfer size.
\item What does this imply for the choice between many small and fewer large
transfers?
\end{enumerate}

\medskip\noindent\textbf{Part IV --- cost-effectiveness.} Administration costs 10~\% of the transfers.
A training programme costs 3{,}000 dalens a household and raises the share above the line by 8 points.
\begin{enumerate}[resume]
\item Find the total cost of the transfers.
\item Find the number of extra households above the line due to the transfers.
\item Find the cost per household lifted above the line.
\item Find the training programme's cost per household lifted.
\item Name two reasons the comparison could mislead.
\item State the result: the effect of the transfer on consumption after two years, with
its interval, and the cost per household lifted above the line.
\end{enumerate}
\end{problem}
